\chapter{Teoría de grupos aplicada}\label{ch:b3:group-theory-applied}

El dióxido de carbono representa unos cientos de moléculas por cada millón en el
aire, y el nitrógeno y el oxígeno casi todo el resto; y, sin embargo, es el dióxido de carbono, y no
el nitrógeno ni el oxígeno, el que absorbe la luz infrarroja que el suelo caliente envía al
espacio. Una molécula solo absorbe luz infrarroja mediante una vibración que
cambia su momento dipolar, y si una vibración lo hace o no lo decide únicamente la
simetría: la única vibración del nitrógeno no puede; dos de las cuatro del dióxido de carbono
sí. Este capítulo convierte las \omterm{def:b3:point-groups:irrep}{tablas de caracteres} del
\cref{ch:b3:point-groups} en herramientas: reduce representaciones, construye
orbitales adaptados a la simetría, cuenta y etiqueta las vibraciones y deduce las
reglas de selección de las espectroscopias infrarroja y Raman.

\begin{recall}
El \cref{ch:b3:point-groups} definió las \omterm{def:b3:point-groups:operation}{operaciones de simetría}, los \omterm{def:b3:point-groups:point-group}{grupos puntuales},
las clases, las representaciones, los \omterm{def:b3:point-groups:representation}{caracteres} y las \omterm{def:b3:point-groups:irrep}{tablas de caracteres}. El volumen del
segundo año construyó los orbitales moleculares de \ce{H2O}, \ce{NH3} y \ce{CH4} a partir de
orbitales de fragmento y de combinaciones adaptadas a la simetría de los orbitales $1s$ de los
hidrógenos, con etiquetas que aquí se deducen. El volumen del primer año
leyó espectros infrarrojos en números de onda y definió la polarizabilidad de una
molécula.
\end{recall}

\section{Reducción de una representación}

Todo conjunto de funciones o de vectores asociado a una molécula lleva una
representación de su \omterm{def:b3:point-groups:point-group}{grupo puntual}. La mayoría son reducibles.

\begin{definition}[Representación reducible, representación totalmente simétrica]\label{def:b3:group-theory-applied:reducible}
Una \emph{representación reducible}\index{representacion reducible@representación reducible} es aquella
cuyas matrices pueden llevarse todas, mediante un mismo cambio de base, a la misma
forma diagonal por bloques; es entonces una suma de \omterm{def:b3:point-groups:irrep}{representaciones irreducibles},
que se escribe $\Gamma = \sum_in_i\Gamma_i$. La \omterm{def:b3:point-groups:irrep}{representación irreducible} cuyos
\omterm{def:b3:point-groups:representation}{caracteres} valen todos 1 es la \emph{representación totalmente
simétrica}\index{representacion totalmente simetrica@representación totalmente simétrica} ($A_1$, $A_{1g}$,
$A_1'$\dots).
\end{definition}

\begin{theorem}[Gran teorema de ortogonalidad]\label{thm:b3:group-theory-applied:great-orthogonality}
Para dos \omterm{def:b3:point-groups:irrep}{representaciones irreducibles} $\Gamma_i$, $\Gamma_j$, de dimensiones $d_i$,
$d_j$, de un grupo de orden $h$, escritas con matrices unitarias,
\[
\sum_R D_i(R)_{mn}^*\,D_j(R)_{m'n'} = \frac{h}{d_i}\,\delta_{ij}\,\delta_{mm'}\,\delta_{nn'}.
\]
\end{theorem}
\admitted

La demostración, mediante los lemas de Schur, corresponde a cursos más avanzados; lo que
usa la química son sus consecuencias para los \omterm{def:b3:point-groups:representation}{caracteres}, y todas las \omterm{def:b3:point-groups:irrep}{tablas de caracteres}
de este libro se comprueban con ellas.

\begin{corollary}[Ortogonalidad de los caracteres]\label{cor:b3:group-theory-applied:characters}
Con $g_c$ el número de operaciones de la clase $c$,
\[
\sum_cg_c\,\chi_i(c)^*\chi_j(c) = h\,\delta_{ij},
\]
y el número de \omterm{def:b3:point-groups:irrep}{representaciones irreducibles} es igual al número de clases,
con $\sum_id_i^2 = h$.
\end{corollary}

\begin{proof}
Se toman $m = n$ y $m' = n'$ en el teorema y se suma sobre $m$ y $m'$: el primer
miembro se convierte en $\sum_R\chi_i(R)^*\chi_j(R)$, el segundo en $\frac{h}{d_i}\delta_{ij}
\sum_{m}\delta_{mm} = h\delta_{ij}$; se agrupan las operaciones por clases. Las filas de
la tabla son, pues, vectores ortogonales en un espacio cuya dimensión es el
número de clases, de modo que hay como mucho tantas \omterm{def:b3:point-groups:irrep}{representaciones
irreducibles} como clases; la igualdad y $\sum d_i^2 = h$ se obtienen aplicando
el teorema a la representación regular (ejercicio 10).
\end{proof}

\begin{theorem}[La fórmula de reducción]\label{thm:b3:group-theory-applied:reduction}
Una representación $\Gamma$ de \omterm{def:b3:point-groups:representation}{caracteres} $\chi(c)$ contiene la \omterm{def:b3:point-groups:irrep}{representación
irreducible} $\Gamma_i$ exactamente
\[
n_i = \frac1h\sum_cg_c\,\chi(c)\,\chi_i(c)^*
\]
veces. Es la \emph{fórmula de reducción}\index{formula de reduccion@fórmula de reducción}.
\end{theorem}

\begin{proof}
Como $\Gamma = \sum_jn_j\Gamma_j$, sus \omterm{def:b3:point-groups:representation}{caracteres} son $\chi(c) = \sum_jn_j\chi_j(c)$.
Se multiplica por $g_c\chi_i(c)^*$, se suma sobre las clases y se usa el corolario:
$\sum_cg_c\chi(c)\chi_i(c)^* = \sum_jn_jh\delta_{ij} = hn_i$.
\end{proof}

\begin{example}[Los orbitales de los hidrógenos del amoniaco]\label{ex:b3:group-theory-applied:nh3}
Se toman como base los tres orbitales $1s$ de los hidrógenos de \ce{NH3}. $E$
deja los tres en su sitio ($\chi = 3$), $C_3$ mueve los tres ($\chi = 0$), un
$\sigma_v$ deja uno en su sitio e intercambia dos ($\chi = 1$). Con la tabla de $C_{3v}$:
$n_{A_1} = \frac16(3 + 0 + 3\times1) = 1$, $n_{A_2} = \frac16(3 + 0 - 3) = 0$,
$n_E = \frac16(6 + 0 + 0) = 1$. Así, $\Gamma_{\mathrm H} = A_1 + E$: los tres
hidrógenos dan una combinación totalmente simétrica y un par degenerado.
\end{example}

\begin{method}[Reducir una representación]\label{met:b3:group-theory-applied:reduce}
\begin{enumerate}
  \item Para cada clase, encontrar el \omterm{def:b3:point-groups:representation}{carácter}: para un conjunto de funciones centradas en los
        átomos, el número de funciones que quedan en su sitio, cada una contada con
        el signo que adquiere.
  \item Aplicar la \omterm{thm:b3:group-theory-applied:reduction}{fórmula de reducción} con la fila de cada \omterm{def:b3:point-groups:irrep}{representación
        irreducible}, ponderando cada clase por su tamaño.
  \item Comprobar: los resultados son enteros no negativos, y $\sum_in_id_i$
        es igual al \omterm{def:b3:point-groups:representation}{carácter} bajo $E$ (el tamaño de la base).
\end{enumerate}
\end{method}

\section{Orbitales adaptados a la simetría}

\begin{definition}[Operador de proyección]\label{def:b3:group-theory-applied:projection}
El \emph{operador de proyección}\index{operador de proyeccion@operador de proyección} sobre la
\omterm{def:b3:point-groups:irrep}{representación irreducible} $\Gamma_i$ es
\[
\hat P_i = \frac{d_i}{h}\sum_R\chi_i(R)^*\,\hat R,
\]
donde $\hat R$ aplica la operación $R$ a una función.
\end{definition}

\begin{proposition}[Qué hace el operador de proyección]\label{prop:b3:group-theory-applied:projection}
Aplicado a cualquier función $f$, $\hat P_if$ es o bien nula o bien una función que
se transforma como $\Gamma_i$ (una combinación adaptada a la simetría); para una
$\Gamma_i$ unidimensional, no cambia, salvo por el signo $\chi_i(S)$, bajo cada
operación $S$.
\end{proposition}

\begin{proof}[Demostración parcial]
Para $d_i = 1$: $\hat S\hat P_if = \frac1h\sum_R\chi_i(R)\hat S\hat Rf$. Sea $R' = SR$,
que recorre todo el grupo cuando lo recorre $R$; $\chi_i(R) = \chi_i(S^{-1}R') =
\chi_i(S)\chi_i(R')$ para una representación unidimensional (los \omterm{def:b3:point-groups:representation}{caracteres} son
las propias matrices). Por tanto, $\hat S\hat P_if = \chi_i(S)\hat P_if$. El caso
general usa el gran teorema de ortogonalidad y se admite.
\end{proof}

\begin{method}[Construir combinaciones adaptadas a la simetría]\label{met:b3:group-theory-applied:salc}
\begin{enumerate}
  \item Elegir una función del conjunto, $h_1$, y tabular la función a la que
        la envía cada operación.
  \item Para cada \omterm{def:b3:point-groups:irrep}{representación irreducible}, multiplicar cada imagen por el
        \omterm{def:b3:point-groups:representation}{carácter} de la operación, sumar y normalizar (despreciando el solapamiento
        entre las funciones atómicas).
  \item Para una representación degenerada, proyectar una segunda función para obtener
        un segundo miembro y ortogonalizar después.
\end{enumerate}
\end{method}

\begin{example}[Las combinaciones del amoniaco y del agua]\label{ex:b3:group-theory-applied:salcs}
En \ce{NH3}, $\hat P_{A_1}h_1 \propto h_1 + h_2 + h_3$ (todos los \omterm{def:b3:point-groups:representation}{caracteres} valen 1). Para $E$
(\omterm{def:b3:point-groups:representation}{caracteres} $2, -1, 0$), $\hat P_Eh_1 \propto 2h_1 - h_2 - h_3$; proyectar $h_2$ y
ortogonalizar da la pareja $h_2 - h_3$. Normalizadas: $a_1 =
\frac{1}{\sqrt3}(h_1 + h_2 + h_3)$, $e = \frac{1}{\sqrt6}(2h_1 - h_2 - h_3)$,
$\frac{1}{\sqrt2}(h_2 - h_3)$. En \ce{H2O}, situada en el plano $xz$, $h_1 + h_2$ es
$a_1$ y $h_1 - h_2$ es $b_1$ (cambia de signo con $C_2$ y con $\sigma_v'(yz)$,
que intercambian los hidrógenos). Con la molécula en el plano $yz$, como prefieren algunos
libros, las etiquetas $b_1$ y $b_2$ se intercambian en todas partes.
\end{example}

\begin{omfigure}
\begin{tikzpicture}[font=\small]
  \foreach \x/\lab/\ca/\cb/\cc in {0/{$a_1$}/1/1/1, 4/{$e$}/2/-1/-1, 8/{$e$}/0/1/-1}{
    \begin{scope}[xshift=\x cm]
      \foreach \a/\c in {90/\ca,210/\cb,330/\cc}{
        \pgfmathsetmacro{\r}{0.18+0.13*abs(\c)}
        \ifdim \c pt > 0pt \fill[omDef!55, draw=omDef] (\a:1.0) circle (\r);
        \else \ifdim \c pt < 0pt \fill[omThm!25, draw=omThm] (\a:1.0) circle (\r);
        \else \draw[gray, densely dotted] (\a:1.0) circle (0.12); \fi\fi
      }
      \node[aN, minimum size=4.6mm] at (0,0) {};
      \node[anchor=north] at (0,-1.45) {\lab};
    \end{scope}}
  \node[anchor=north, font=\footnotesize] at (0,-1.85) {$h_1 + h_2 + h_3$};
  \node[anchor=north, font=\footnotesize] at (4,-1.85) {$2h_1 - h_2 - h_3$};
  \node[anchor=north, font=\footnotesize] at (8,-1.85) {$h_2 - h_3$};
\end{tikzpicture}

{\small Las combinaciones adaptadas a la simetría de los tres orbitales $1s$ de los hidrógenos
del amoniaco, vistas según el eje $C_3$ ($h_1$ arriba). Azul: coeficiente
positivo; rojo: negativo; el tamaño de cada disco crece con el
coeficiente; punteado: nulo.}
\end{omfigure}

\begin{example}[Seis ligandos en torno a un metal]\label{ex:b3:group-theory-applied:ml6}
Para seis orbitales $\sigma$ de ligandos que apuntan hacia un metal desde los vértices de un
octaedro, los \omterm{def:b3:point-groups:representation}{caracteres} sobre las clases de $O_h$ ($E$, $8C_3$, $6C_2$, $6C_4$,
$3C_2$, $i$, $6S_4$, $8S_6$, $3\sigma_h$, $6\sigma_d$) son $6, 0, 0, 2, 2, 0, 0, 0, 4,
2$, y la reducción da $\Gamma_\sigma = A_{1g} + E_g + T_{1u}$. Los orbitales $s$
($a_{1g}$), $d_{z^2}$ y $d_{x^2-y^2}$ ($e_g$) y $p$ ($t_{1u}$) del metal encuentran
pareja entre los ligandos; sus $d_{xy}$, $d_{xz}$, $d_{yz}$ ($t_{2g}$) no la encuentran y quedan no enlazantes
en $\sigma$ ---el origen del desdoblamiento $t_{2g}$/$e_g$ del volumen del segundo año,
desarrollado en los \cref{ch:b3:organometallic-bonding,ch:b3:complex-spectra-magnetism}.
\end{example}

\begin{omfigure}
\begin{tikzpicture}[font=\small, yscale=0.95]
  % O orbitals (left), MOs (centre), H combinations (right); schematic energies
  \foreach \y/\l in {0/{$2s\ (a_1)$}}{\draw[omstate] (0,\y) -- (1.2,\y); \node[anchor=east] at (-0.05,\y) {\l};}
  \draw[omstate] (0,3.0) -- (0.36,3.0) (0.42,3.0) -- (0.78,3.0) (0.84,3.0) -- (1.2,3.0);
  \node[anchor=east] at (-0.05,3.0) {$2p\ (b_1, a_1, b_2)$};
  \draw[omstate] (6.8,3.9) -- (8.0,3.9); \node[anchor=west] at (8.05,3.9) {$h_1 - h_2\ (b_1)$};
  \draw[omstate] (6.8,3.6) -- (8.0,3.6); \node[anchor=west] at (8.05,3.5) {$h_1 + h_2\ (a_1)$};
  \foreach \y/\l in {-0.6/{$2a_1$}, 1.6/{$1b_1$}, 2.4/{$3a_1$}, 3.0/{$1b_2$}, 5.6/{$4a_1$}, 6.2/{$2b_1$}}{
    \draw[omstate] (3.4,\y) -- (4.6,\y); \node[anchor=west] at (4.65,\y) {\l};}
  \foreach \y in {-0.6,1.6,2.4,3.0}{\node[font=\scriptsize] at (4.0,\y+0.16) {$\uparrow\downarrow$};}
  \draw[densely dotted, gray] (1.2,0) -- (3.4,-0.6) (1.2,3.0) -- (3.4,1.6) (1.2,3.0) -- (3.4,2.4) (1.2,3.0) -- (3.4,3.0)
     (6.8,3.6) -- (4.6,-0.6) (6.8,3.9) -- (4.6,1.6) (6.8,3.6) -- (4.6,2.4) (1.2,3.0) -- (3.4,6.2) (6.8,3.9) -- (4.6,6.2) (6.8,3.6) -- (4.6,5.6);
  \node[anchor=north] at (0.6,-1.0) {O};
  \node[anchor=north] at (4.0,-1.0) {\ce{H2O}};
  \node[anchor=north] at (7.4,-1.0) {2 H};
  \draw[->] (-2.3,-0.9) -- (-2.3,6.4) node[above] {$E$};
\end{tikzpicture}

{\small Orbitales moleculares de valencia del agua (energías esquemáticas, molécula en
el plano $xz$). Solo se mezclan orbitales de la misma simetría: la combinación $b_1$ de los
hidrógenos con $2p_x$, la $a_1$ con $2s$ y $2p_z$; $2p_y$ ($b_2$) no tiene
pareja y queda como par no enlazante, el orbital ocupado más alto.}
\end{omfigure}

\section{Modos de vibración}

\begin{definition}[Modo normal]\label{def:b3:group-theory-applied:normal-mode}
Un \emph{modo normal}\index{modo normal} de una molécula es una vibración
colectiva en la que todos los átomos oscilan a la misma frecuencia y en fase,
según el vector propio de la hessiana ponderada por la masa
(\cref{prop:b3:computational-chemistry:hessian}). Cada modo normal
se transforma según una \omterm{def:b3:point-groups:irrep}{representación irreducible} del \omterm{def:b3:point-groups:point-group}{grupo puntual}.
\end{definition}

\begin{proposition}[La representación de todos los movimientos]\label{prop:b3:group-theory-applied:gamma-3n}
Se asocian tres vectores de desplazamiento $x$, $y$, $z$ a cada átomo. Para una operación
$R$, el \omterm{def:b3:point-groups:representation}{carácter} de esta representación de dimensión $3N$ es
\[
\chi_{3N}(R) = (\text{número de átomos que quedan en su sitio}) \times \chi_{xyz}(R),
\]
con $\chi_{xyz} = 1 + 2\cos\theta$ para una rotación de ángulo $\theta$ y $-1 + 2\cos\theta$
para una rotación impropia (es decir, 3 para $E$, $-1$ para $C_2$, 0 para $C_3$, 1 para $\sigma$,
$-3$ para $i$).
\end{proposition}

\begin{proof}
Un átomo desplazado a otra posición lleva sus vectores fuera de la diagonal de
la matriz: no contribuye a la traza. Un átomo que queda en su sitio tiene
sus tres vectores transformados por la matriz $3\times3$ de $R$, cuya traza se
calcula en un referencial con $z$ según el eje: $\cos\theta + \cos\theta + 1$ para una
rotación, y el último término pasa a $-1$ para una rotación impropia.
\end{proof}

\begin{proposition}[Recuento de las vibraciones]\label{prop:b3:group-theory-applied:mode-count}
$\Gamma_{\mathrm{vib}} = \Gamma_{3N} - \Gamma_{\mathrm{trans}} - \Gamma_{\mathrm{rot}}$, donde
$\Gamma_{\mathrm{trans}}$ es la representación de $(x, y, z)$ y $\Gamma_{\mathrm{rot}}$ la
de $(R_x, R_y, R_z)$; una molécula no lineal tiene $3N - 6$ vibraciones, una lineal
$3N - 5$.
\end{proposition}

\begin{proof}
Los $3N$ desplazamientos describen todos los movimientos: tres traslaciones del
centro de masas, tres rotaciones (dos para una molécula lineal, pues girar
en torno a su propio eje no mueve ningún núcleo) y las vibraciones. Estos subespacios no
se mezclan con las operaciones, de modo que sus representaciones se suman.
\end{proof}

\begin{example}[El agua]\label{ex:b3:group-theory-applied:water}
% ledger: vib:H2O.sym, vib:H2O.bend, vib:H2O.asym
\ce{H2O} en $C_{2v}$, molécula en el plano $xz$. Átomos que quedan en su sitio: 3, 1, 3, 1;
$\chi_{xyz}$: $3, -1, 1, 1$; luego $\chi_{3N} = 9, -1, 3, 1$. La reducción da $3A_1 +
A_2 + 3B_1 + 2B_2$. Una vez restadas las traslaciones $A_1 + B_1 + B_2$ y las rotaciones $A_2 + B_1 + B_2$,
queda $\Gamma_{\mathrm{vib}} = 2A_1 + B_1$. Los dos modos $a_1$ son la tensión
simétrica (\qty{3657}{cm^{-1}}) y la flexión (\qty{1595}{cm^{-1}}); el modo $b_1$ es
la tensión antisimétrica (\qty{3756}{cm^{-1}}).
\end{example}

\begin{omfigure}
\begin{tikzpicture}[font=\small, >=Stealth]
  \foreach \x/\name/\lab in {0/sym/{$\nu_1$, $a_1$, tensión simétrica}, 4.3/bend/{$\nu_2$, $a_1$, flexión}, 8.6/asym/{$\nu_3$, $b_1$, tensión antisimétrica}}{
    \begin{scope}[xshift=\x cm]
      \node[aO] (o\name) at (0,0.35) {};
      \node[aH] (a\name) at (-0.8,-0.25) {};
      \node[aH] (b\name) at (0.8,-0.25) {};
      \begin{scope}[on background layer]\draw[bond] (o\name.center) -- (a\name.center) (o\name.center) -- (b\name.center);\end{scope}
      \node[anchor=north, align=center, font=\footnotesize] at (0,-0.9) {\lab};
    \end{scope}}
  % symmetric stretch: H outward along the bonds, O up a little
  \draw[->, omThm, thick] (-0.95,-0.36) -- (-1.42,-0.71);
  \draw[->, omThm, thick] (0.95,-0.36) -- (1.42,-0.71);
  \draw[->, omThm, thick] (0,0.74) -- (0,0.98);
  % bend: each H moves perpendicular to its bond, closing the angle; O up
  \draw[->, omThm, thick] (4.3-0.92,-0.36) -- (4.3-0.62,-0.76);
  \draw[->, omThm, thick] (4.3+0.92,-0.36) -- (4.3+0.62,-0.76);
  \draw[->, omThm, thick] (4.3,0.74) -- (4.3,0.98);
  % antisymmetric stretch: one H out, the other in, O sideways
  \draw[->, omThm, thick] (8.6-0.95,-0.36) -- (8.6-1.42,-0.71);
  \draw[->, omThm, thick] (8.6+1.02,0.0) -- (8.6+0.70,0.24);
  \draw[->, omThm, thick] (8.6-0.05,0.80) -- (8.6+0.25,0.80);
\end{tikzpicture}

{\small Los tres \omterm{def:b3:group-theory-applied:normal-mode}{modos normales} del agua (flechas: desplazamientos, no a
escala). Los dos modos $a_1$ conservan toda la simetría de la molécula; el modo $b_1$
cambia de signo con $C_2$.}
\end{omfigure}

\begin{example}[Cuatro moléculas más]\label{ex:b3:group-theory-applied:table}
El mismo procedimiento (calculado y comprobado en los datos de las figuras de este
capítulo) da:
\begin{center}\small
\begin{tabular}{llll}
\hline
molécula & grupo & $\Gamma_{\mathrm{vib}}$ & modos \\ \hline
\ce{NH3} & $C_{3v}$ & $2A_1 + 2E$ & 6 \\
\ce{CH4} & $T_d$ & $A_1 + E + 2T_2$ & 9 \\
\ce{CO2} & $D_{\infty h}$ (vía $D_{2h}$) & $\Sigma_g^+ + \Sigma_u^+ + \Pi_u$ ($A_g + B_{1u} + B_{2u} + B_{3u}$) & 4 \\
\ce{BF3} & $D_{3h}$ & $A_1' + 2E' + A_2''$ & 6 \\
\ce{XeF4} & $D_{4h}$ & $A_{1g} + B_{1g} + B_{2g} + A_{2u} + B_{2u} + 2E_u$ & 9 \\
\ce{SF6} & $O_h$ & $A_{1g} + E_g + T_{2g} + 2T_{1u} + T_{2u}$ & 15 \\ \hline
\end{tabular}
\end{center}
Para una molécula lineal, el grupo infinito se sustituye por su \omterm{def:b3:point-groups:class}{subgrupo} $D_{2h}$,
en el que la flexión degenerada $\Pi_u$ se desdobla en $B_{2u} + B_{3u}$.
\end{example}

\section{Reglas de selección}

\begin{definition}[Producto directo]\label{def:b3:group-theory-applied:direct-product}
El \emph{producto directo}\index{producto directo} $\Gamma_i\otimes\Gamma_j$ de dos
representaciones es la representación que llevan los productos de sus
funciones de base; sus \omterm{def:b3:point-groups:representation}{caracteres} son los productos $\chi_i(R)\chi_j(R)$.
\end{definition}

\begin{theorem}[Integrales nulas]\label{thm:b3:group-theory-applied:vanishing-integral}
Una integral $\int f_1f_2f_3\,\dd\tau$ extendida a todo el espacio solo puede ser no nula si el
\omterm{def:b3:group-theory-applied:direct-product}{producto directo} $\Gamma_1\otimes\Gamma_2\otimes\Gamma_3$ contiene la \omterm{def:b3:group-theory-applied:reducible}{representación
totalmente simétrica}.
\end{theorem}

\begin{proof}
Una \omterm{def:b3:point-groups:operation}{operación de simetría} solo vuelve a etiquetar los puntos del espacio, de modo que el valor de la
integral no cambia cuando el integrando se transforma por cualquier $R$. Se promedia
el integrando transformado sobre el grupo: el promedio de las componentes del
integrando que pertenecen a cada \omterm{def:b3:point-groups:irrep}{representación irreducible} es
$\frac1h\sum_R\hat R(\cdot)$, que es la proyección sobre la \omterm{def:b3:group-theory-applied:reducible}{representación totalmente
simétrica} (\cref{def:b3:group-theory-applied:projection} con todos los
$\chi = 1$). Todas las demás componentes tienen promedio nulo; si el producto no contiene ninguna
parte totalmente simétrica, la integral es nula.
\end{proof}

\begin{definition}[Activo en IR, activo en Raman]\label{def:b3:group-theory-applied:activity}
Una vibración fundamental es \emph{activa en IR}\index{activo en IR} si puede absorber
luz infrarroja, y \emph{activa en Raman}\index{activo en Raman} si aparece en el
espectro Raman (\cref{ch:b3:rovibrational-spectroscopy}).
\end{definition}

\begin{corollary}[Reglas de selección infrarroja y Raman]\label{cor:b3:group-theory-applied:ir-raman}
Una fundamental (del nivel fundamental, totalmente simétrico, a un cuanto
de un modo de simetría $\Gamma$) solo es \omterm{def:b3:group-theory-applied:activity}{activa en IR} si $\Gamma$ es la
representación de $x$, $y$ o $z$, y solo es \omterm{def:b3:group-theory-applied:activity}{activa en Raman} si $\Gamma$ es la
representación de una función cuadrática ($x^2$, $xy$, \dots).
\end{corollary}

\begin{proof}
La intensidad de absorción hace intervenir $\int\psi_0\,\mu_k\,\psi_1\dd\tau$ con las componentes
del dipolo $\mu_k$, que se transforman como $x$, $y$, $z$; $\psi_0$ es totalmente simétrica
y $\psi_1$ se transforma como $\Gamma$. Por el teorema, la integral solo puede ser no nula
si $\Gamma\otimes\Gamma_{\mu_k}$ contiene $A_1$, lo que solo ocurre si $\Gamma =
\Gamma_{\mu_k}$ (el producto de dos \omterm{def:b3:point-groups:irrep}{representaciones irreducibles} solo contiene la
totalmente simétrica si son iguales, para \omterm{def:b3:point-groups:representation}{caracteres} reales). La dispersión Raman
hace intervenir las componentes de la polarizabilidad $\alpha_{kl}$, que se
transforman como los productos $kl$.
\end{proof}

\begin{proposition}[Regla de exclusión mutua]\label{prop:b3:group-theory-applied:mutual-exclusion}
En una molécula con \omterm{def:b3:point-groups:inversion}{centro de inversión}, ninguna fundamental es a la vez \omterm{def:b3:group-theory-applied:activity}{activa en IR} y en
Raman: es la \emph{regla de exclusión mutua}\index{regla de exclusion mutua@regla de exclusión mutua}.
\end{proposition}

\begin{proof}
Con un \omterm{def:b3:point-groups:inversion}{centro de inversión}, toda \omterm{def:b3:point-groups:irrep}{representación irreducible} es $g$ o $u$. Las
coordenadas $x$, $y$, $z$ cambian de signo con $i$ (son $u$); las funciones
cuadráticas no ($g$). Un modo solo es \omterm{def:b3:group-theory-applied:activity}{activo en IR} si es $u$, y solo \omterm{def:b3:group-theory-applied:activity}{activo en Raman} si es
$g$.
\end{proof}

\begin{example}[El dióxido de carbono y el efecto invernadero]\label{ex:b3:group-theory-applied:co2}
% ledger: vib:CO2.sym, vib:CO2.bend, vib:CO2.asym, ir:CO2.asym.int, ir:CO2.bend.int
\ce{CO2} tiene un \omterm{def:b3:point-groups:inversion}{centro de inversión}. Su tensión simétrica ($\Sigma_g^+$,
\qty{1333}{cm^{-1}}) solo es \omterm{def:b3:group-theory-applied:activity}{activa en Raman}; la tensión antisimétrica ($\Sigma_u^+$,
\qty{2349}{cm^{-1}}) y la flexión ($\Pi_u$, \qty{667}{cm^{-1}}) solo son \omterm{def:b3:group-theory-applied:activity}{activas en IR}.
La flexión absorbe cerca de \qty{15}{\micro m}, cerca del máximo de la emisión
térmica de la Tierra: por eso el dióxido de carbono es un gas de efecto invernadero. \ce{N2} y
\ce{O2} tienen una sola vibración, $\Sigma_g^+$: inactiva en IR.
\end{example}

\begin{omfigure}
\begin{tikzpicture}
\begin{axis}[omir, width=0.92\linewidth, height=4.4cm, xmin=400, xmax=2600,
  ymin=0, ymax=1.3, ytick=\empty, xlabel={número de onda / cm$^{-1}$},
  ylabel={señal}, legend cell align=left,
  legend style={font=\small, at={(0.98,0.95)}, anchor=north east, draw=none}]
\addplot[omDef, thick] table[x=nu, y=ir] {figdata/out/bachelor-3-co2-spectra.dat};
\addlegendentry{absorción infrarroja}
\addplot[omThm, thick, dashed] table[x=nu, y=raman] {figdata/out/bachelor-3-co2-spectra.dat};
\addlegendentry{dispersión Raman}
\node[font=\small, anchor=south] at (axis cs:2349,1.02) {$\Sigma_u^+$};
\node[font=\small, anchor=south] at (axis cs:667,0.1) {$\Pi_u$};
\node[font=\small, anchor=south, omThm] at (axis cs:1333,1.02) {$\Sigma_g^+$};
\end{axis}
\end{tikzpicture}

{\small Espectros sintéticos de \ce{CO2} gaseoso (posiciones de las bandas e intensidades
infrarrojas relativas tomadas de los valores medidos; formas esquemáticas, sin
estructura de rotación). Exclusión mutua: ninguna banda aparece en ambos.}
\end{omfigure}

\begin{method}[Predecir un espectro de vibración]\label{met:b3:group-theory-applied:vibrations}
\begin{enumerate}
  \item Encontrar el \omterm{def:b3:point-groups:point-group}{grupo puntual}; calcular $\chi_{3N}$ a partir de los átomos que quedan en su sitio.
  \item Reducir y restar traslaciones y rotaciones: $\Gamma_{\mathrm{vib}}$.
  \item Con las columnas de la derecha de la tabla, marcar cada modo como \omterm{def:b3:group-theory-applied:activity}{activo en IR}
        ($x$, $y$, $z$) y/o \omterm{def:b3:group-theory-applied:activity}{activo en Raman} (cuadrático); un modo que no es ni lo uno ni lo otro
        es silencioso.
  \item Contar las bandas: una por modo activo; un modo degenerado da una sola banda.
\end{enumerate}
\end{method}

\begin{method}[Contar las tensiones CO]\label{met:b3:group-theory-applied:co-count}
En un carbonilo metálico, se toman como base los $n$ vectores de enlace C--O (un átomo que queda en
su sitio cuenta 1, nada más): se reduce para obtener $\Gamma_{\mathrm{CO}}$ y después se aplican
las reglas IR y Raman. El número de bandas intensas en torno a
\qtyrange{1850}{2150}{cm^{-1}} identifica el isómero.
\end{method}

\begin{example}[Cis y trans]\label{ex:b3:group-theory-applied:cis-trans}
cis-\ce{ML4(CO)2}, $C_{2v}$: \omterm{def:b3:point-groups:representation}{caracteres} $2, 0, 2, 0$, $\Gamma_{\mathrm{CO}} = A_1 + B_1$,
dos bandas IR. trans-\ce{ML4(CO)2}, $D_{4h}$: $\Gamma_{\mathrm{CO}} = A_{1g} + A_{2u}$, una
banda IR ($A_{2u}$) y una banda Raman ($A_{1g}$). Una sola banda CO en el
infrarrojo indica el isómero trans.
\end{example}

\section{Aplicaciones a las transiciones electrónicas}

El mismo teorema se aplica a las transiciones electrónicas: una transición electrónica
entre los estados $\Psi_1$ y $\Psi_2$ solo está permitida si $\Gamma_1\otimes\Gamma_2$
contiene la representación de $x$, $y$ o $z$, y su luz está polarizada según
ese eje.

\begin{example}[Transiciones del agua y del formaldehído]\label{ex:b3:group-theory-applied:electronic}
En el agua ($C_{2v}$, plano $xz$), promover un electrón de $1b_2$ (el par no enlazante)
a $4a_1$ da un estado excitado de simetría $B_2\otimes A_1 = B_2$, que es $y$:
permitida, polarizada perpendicularmente al plano molecular. En el formaldehído
($C_{2v}$), la transición $n\to\pi^*$ va de un par no enlazante del oxígeno contenido en el plano ($b_1$)
al orbital $\pi^*$, construido con orbitales $p$ perpendiculares al plano
($b_2$; molécula en el plano $xz$, C=O según $z$): $B_1\otimes B_2 = A_2$,
que no es ninguno de $x$, $y$, $z$: prohibida por simetría, y por eso su banda
(\cref{ch:b3:electronic-spectroscopy}) es tan débil.
\end{example}

\begin{history}[Wigner, Bethe y la simetría de las moléculas]
La teoría de grupos entró en la mecánica cuántica con el trabajo de Eugene Wigner sobre los espectros
atómicos (1927--31) y el artículo de Hans Bethe sobre el desdoblamiento de los términos atómicos en los
cristales (1929). E. Bright Wilson la aplicó a las vibraciones moleculares en
1934; Robert Mulliken dio las etiquetas que aún se usan para orbitales y estados.
\end{history}

\begin{inthelab}[Infrarrojo y Raman, uno junto al otro]
Un espectrómetro infrarrojo mide la luz que absorbe una muestra; un espectrómetro
Raman la ilumina con un láser y analiza la débil luz dispersada
a longitudes de onda desplazadas. Registrar ambos espectros del mismo compuesto es la
prueba clásica de la existencia de un centro de simetría: si ninguna banda coincide, la molécula
es probablemente centrosimétrica. Los láseres Raman son de clase 3B o 4: el haz está
confinado y los \omterm{def:b3:quantum-model-systems:operator}{operadores} llevan gafas homologadas para su longitud de onda.
\end{inthelab}

\section{Ejercicios}

\begin{exercise}[$\star$]\label{exo:b3:group-theory-applied:1}
Redúzcase la representación de $C_{2v}$ de \omterm{def:b3:point-groups:representation}{caracteres} $(5, -1, 3, 1)$ sobre $(E, C_2,
\sigma_v(xz), \sigma_v'(yz))$. Compruébese la dimensión.
\end{exercise}

\begin{exercise}[$\star$]\label{exo:b3:group-theory-applied:2}
\ce{SO2} es angular, como el agua. Dese $\Gamma_{\mathrm{vib}}$ e indíquese qué modos son \omterm{def:b3:group-theory-applied:activity}{activos en IR}
y en Raman.
\end{exercise}

\begin{exercise}[$\star$]\label{exo:b3:group-theory-applied:3}
En $C_{2v}$, calcúlense los \omterm{def:b3:group-theory-applied:direct-product}{productos directos} $B_1\otimes B_2$ y $A_2\otimes B_1$, y
$E\otimes E$ en $C_{3v}$ (redúzcase este último).
\end{exercise}

\begin{exercise}[$\star$]\label{exo:b3:group-theory-applied:4}
% ledger: vib:CH4.nu1, vib:CH4.nu2, vib:CH4.nu3, vib:CH4.nu4
\ce{CH4} tiene fundamentales a 2917 ($a_1$), 1534 ($e$), 3019 y 1306 (ambas $t_2$)
\unit{cm^{-1}}. ¿Cuáles aparecen en el espectro infrarrojo y cuáles en el espectro
Raman?
\end{exercise}

\begin{exercise}[$\star\star$]\label{exo:b3:group-theory-applied:5}
Para \ce{BF3}, dedúzcase $\Gamma_{\mathrm{vib}} = A_1' + 2E' + A_2''$ a partir de los átomos que quedan en
su sitio, y cuéntense después las bandas IR y Raman. ¿Qué modo se ve en uno solo de
los dos espectros y cuál en ambos?
\end{exercise}

\begin{exercise}[$\star\star$]\label{exo:b3:group-theory-applied:6}
¿Cuántas bandas IR y Raman presenta \ce{SF6}? ¿Qué modo es silencioso? ¿Por qué
ninguna banda aparece en ambos espectros?
\end{exercise}

\begin{exercise}[$\star\star$]\label{exo:b3:group-theory-applied:7}
Aplíquese el \omterm{def:b3:group-theory-applied:projection}{operador de proyección} a $h_2$ para la representación $E$ de $C_{3v}$ y
demuéstrese que ortogonalizar el resultado respecto de $2h_1 - h_2 - h_3$ da $h_2 - h_3$.
\end{exercise}

\begin{exercise}[$\star\star$]\label{exo:b3:group-theory-applied:8}
Predígase el número de bandas IR de CO de fac- y mer-\ce{ML3(CO)3}, y de
\ce{M(CO)6}.
\end{exercise}

\begin{exercise}[$\star\star$]\label{exo:b3:group-theory-applied:9}
Demuéstrese que los seis orbitales $\sigma$ de los ligandos de un complejo octaédrico generan
$A_{1g} + E_g + T_{1u}$, e indíquese con qué orbitales del metal puede combinarse cada uno.
\end{exercise}

\begin{exercise}[$\star\star\star$]\label{exo:b3:group-theory-applied:10}
La representación regular tiene $\chi(E) = h$ y $\chi(R) = 0$ para $R \ne E$. Demuéstrese
con la \omterm{thm:b3:group-theory-applied:reduction}{fórmula de reducción} que contiene cada \omterm{def:b3:point-groups:irrep}{representación irreducible}
$\Gamma_i$ exactamente $d_i$ veces, y dedúzcase $\sum_id_i^2 = h$.
\end{exercise}

\begin{exercise}[$\star\star\star$]\label{exo:b3:group-theory-applied:11}
El acetileno \ce{HC#CH} es lineal. Trabajando en $D_{2h}$, encuéntrese $\Gamma_{\mathrm{vib}}$, reagrúpese
con las etiquetas de $D_{\infty h}$ y predíganse los espectros IR y Raman.
\end{exercise}

\begin{exercise}[$\star\star\star$]\label{exo:b3:group-theory-applied:12}
En el agua (plano $xz$), ¿cuáles de estas promociones monoelectrónicas dan transiciones
permitidas, y con qué polarización: $1b_2 \to 4a_1$, $3a_1 \to 4a_1$,
$1b_2 \to 2b_1$, $1b_1 \to 4a_1$?
\end{exercise}

\section{Problema: ¿cis o trans? Un complejo carbonílico leído en su espectro}

\begin{problem}[{Problema de fin de semana --- grupos puntuales y representaciones de las tensiones CO
de cuatro isómeros, sus bandas infrarrojas y Raman, y la
identificación de un producto a partir de su espectro medido}]
\label{pb:b3:group-theory-applied:1}
% ledger: ct:C2v, ct:C3v, ct:D4h
Un metal M octaédrico lleva ligandos carbonilo y otros ligandos L (cada L
contado como un punto). Se consideran cuatro isómeros: cis- y
trans-\ce{ML4(CO)2}, y fac- y mer-\ce{ML3(CO)3}. Se supone que los ligandos L no
rebajan la simetría más allá de lo que imponen sus posiciones. Una síntesis de
\ce{ML3(CO)3} da un producto cuyo espectro infrarrojo presenta tres bandas
intensas a 2048, 1965 y \qty{1938}{cm^{-1}} (datos del problema); su espectro Raman
presenta tres bandas en posiciones casi idénticas.

\medskip\noindent\textbf{Parte I --- \omterm{def:b3:point-groups:point-group}{Grupos puntuales}.}
\begin{enumerate}
  \item Dibújense los cuatro isómeros sobre un octaedro.
  \item Dese el \omterm{def:b3:point-groups:point-group}{grupo puntual} de cis-\ce{ML4(CO)2}.
  \item El de trans-\ce{ML4(CO)2}.
  \item El de fac-\ce{ML3(CO)3}.
  \item El de mer-\ce{ML3(CO)3}.
  \item ¿Cuáles de los cuatro tienen \omterm{def:b3:point-groups:inversion}{centro de inversión}?
\end{enumerate}

\medskip\noindent\textbf{Parte II --- Las tensiones CO.}
\begin{enumerate}[resume]
  \item Explíquese por qué los vectores de enlace C--O llevan una representación en la que una
        operación contribuye con 1 por cada CO que queda en su sitio.
  \item Encuéntrese $\Gamma_{\mathrm{CO}}$ del isómero cis.
  \item La del isómero trans.
  \item La del isómero fac.
  \item La del isómero mer.
  \item Compruébese cada dimensión con el número de ligandos CO.
\end{enumerate}

\medskip\noindent\textbf{Parte III --- Actividad.}
\begin{enumerate}[resume]
  \item Cuéntense las tensiones CO \omterm{def:b3:group-theory-applied:activity}{activas en IR} de cada isómero.
  \item Cuéntense las \omterm{def:b3:group-theory-applied:activity}{activas en Raman}.
  \item ¿Qué isómero cumple la \omterm{prop:b3:group-theory-applied:mutual-exclusion}{regla de exclusión mutua}?
  \item En el isómero trans, descríbase el modo \omterm{def:b3:group-theory-applied:activity}{activo en IR}: ¿vibran los dos CO
        en fase o en oposición de fase?
  \item En el isómero fac, ¿por qué los tres CO solo dan dos bandas?
  \item ¿Qué indicaría un recuento de 1 banda (IR)?
\end{enumerate}

\medskip\noindent\textbf{Parte IV --- El producto.}
\begin{enumerate}[resume]
  \item ¿Qué isómero indica el espectro infrarrojo del producto?
  \item ¿Es coherente el espectro Raman?
  \item La banda más alta es la tensión en fase de los grupos CO. ¿A qué
        \omterm{def:b3:point-groups:irrep}{representación irreducible} pertenece en ese isómero?
  \item ¿Podría proceder el espectro de una mezcla de los dos isómeros? ¿Qué
        otra medida lo decidiría?
  \item Enúnciese el resultado: el número de tensiones CO \omterm{def:b3:group-theory-applied:activity}{activas en IR} del isómero
        mer, frente al del isómero fac.
\end{enumerate}
\end{problem}
